%% pdflatex ratint.tex
%% bibtex ratint
%% pdflatex ratint.tex
%% pdflatex ratint.tex

%% Upstream typeset this as a plain TeX document with a private
%% "macros.tex" which is not part of the JACAL distribution.  It has
%% been converted to LaTeX.  The page layout, section headings, table
%% of contents and lemma style are configured below so that the result
%% still matches the PDF upstream shipped.

\documentclass[10pt,letterpaper]{article}

%% Plain TeX's page: 1in margins, \vsize of 8.9in, folio below that.
\usepackage[letterpaper,top=1in,left=1in,right=1in,
            textheight=8.9in,footskip=23.4pt]{geometry}
\usepackage{amsmath}
\usepackage{amsthm}

%% Plain TeX's paragraph indentation and body-size footnotes, which
%% the original was set with.
\setlength{\parindent}{20pt}
\renewcommand{\footnotesize}{\normalsize}

%% Plain TeX's display-math spacing (LaTeX's is 2pt tighter, which
%% would repaginate the document).
\AtBeginDocument{%
  \setlength{\abovedisplayskip}     {12pt plus 3pt minus 9pt}%
  \setlength{\belowdisplayskip}     {12pt plus 3pt minus 9pt}%
  \setlength{\abovedisplayshortskip} {0pt plus 3pt}%
  \setlength{\belowdisplayshortskip} {7pt plus 3pt minus 4pt}%
}

%% Sections as "1.<quad>Title" at body size, and a table of contents
%% with slanted entries and dotted leaders, as in the original.
\renewcommand{\thesection}{\arabic{section}.}
\renewcommand{\contentsname}{Table of Contents}
\makeatletter
%% Plain TeX's \beginsection encouraged a page break before a heading
%% by offering the page a third of its height of stretch at a -250
%% penalty, and taking it back again if the break was not used.
%% Without this the pages here run full and the sections all move up.
\newcommand*{\sectionbreak}{%
  \par
  \vskip\z@\@plus.3\textheight
  \penalty-250
  \vskip\z@\@plus-.3\textheight}
\renewcommand{\section}{\sectionbreak\@startsection{section}{1}{\z@}%
  {-\bigskipamount}{\smallskipamount}%
  {\normalfont\normalsize\bfseries}}
\renewcommand*{\l@section}[2]{\@dottedtocline{1}{1em}{1.8em}{\slshape#1}{#2}}
\renewcommand{\@tocrmarg}{\@pnumwidth}
\makeatother

%% Plain TeX's \proclaim: bold run-in heading, slanted body.
\newtheoremstyle{proclaim}{\medskipamount}{\medskipamount}%
  {\slshape}{0pt}{\bfseries}{.}{0.5em}{}
\theoremstyle{proclaim}
\newtheorem{lemma}{Lemma}

%% Three equal columns spanning the text width, for Table 1.
\newlength{\thirdwidth}
\setlength{\thirdwidth}{\dimexpr(\textwidth-4\tabcolsep)/3\relax}

%% This document is dated by its release, not by when it is typeset.
\renewcommand{\today}{30 March 2024}

\DeclareMathOperator{\diff}{d}
\DeclareMathOperator{\W}{W}
\DeclareMathOperator{\atanh}{atanh}
\DeclareMathOperator{\erf}{erf}
\DeclareMathOperator{\fre}{fre}
\newcommand{\coeff}{\mathrm{coeff}\,}
\newcommand{\at}{{@}}

%% Plain TeX packs each page into a \vbox to \vsize; article.cls is
%% \raggedbottom for oneside.  Match the original.
\flushbottom

\begin{document}

{\centering
  \textbf{Canonical Rational Function Integration}\\
  Aubrey G. Jaffer\\
  agj@alum.mit.edu\\
  \today\par}

\section*{Abstract}

{\narrower

  \hspace*{\parindent}%% \section* suppresses the indent; the original had it.
  The derivative of any rational function (ratio of polynomials) is a
  rational function.  An algorithm and decision procedure for finding
  the rational function anti-derivative of a rational function is
  presented.  This algorithm is then extended to derivatives of
  rational functions which include instances of a radical involving
  the integration variable.  It is further extended to derivatives of
  rational functions which include instances of transcendental
  functions defined by separable first-order differential equations.

  \par}

%% \beginsection{Keywords}
%% {\narrower symbolic integration\par}

\tableofcontents

\section{Canonical form}

In symbolic computation, a canonical form has only one representation
for expressions which are equivalent.  A less stringent constraint is
normal form, where all expressions equivalent to 0 are represented
by~0.

Richardson~\cite{10.2307/2271358},
Caviness~\cite{Caviness:1970:CFS:321574.321591}, and
Wang~\cite{10.1145/321850.321856} explore which combinations of number
systems, operations, algebraic, and transcendental functions are
decidable (able to be resolved in finite time by an algorithm) for
questions of an expression having a value greater than 0.  Questions
of ordering are valid only in ordered number systems such as the
reals.  Real analysis is an important domain, but not the only one.

Differential algebra (Ritt~\cite{ritt1966differential}) and Gr\"obner
bases~\cite{alma991010801422205596} are built from a polynomial ring
whose polynomial coefficients are a field.  The canonical form of the
present work is built from the Euclidean domain (which is also a
unique factorization domain and a commutative ring) of polynomials
having integer coefficients.\footnote
%%
{The coefficients being integer instead of rational is of practical
 importance to the speed of calculation.  In order to keep the storage
 for rational numbers under control, polynomial arithmetic algorithms
 typically normalize the ratio of integers using greatest common
 divisor (GCD) after each arithmetic operation, trading time for
 storage.  The storage required for the results of operations on
 integers grows more slowly than for rational numbers.  In the present
 work, GCD operations on integers are postponed until the
 normalization phase.}

In the present work, each equation is in implicit form, a polynomial
equated with zero.  An expression is represented by a polynomial
equation involving the special variable $\at$.  The ratio
$a/(2+3\,b^2)$ is represented:
\begin{equation}
2\,\at+3\,b^2\,\at-a=0\label{at-form}
\end{equation}
%%
A square-root of $a$ is represented by $\at^2-2\,a=0$.
A root of the fifth-order polynomial $x^5+x+2\,a$ is represented by
$\at^5+\at+2\,a=0$.

It is expected that the polynomial arithmetic algorithms will always
combine the coefficients of the sum of terms with identical variable
exponents, and remove monomial terms having 0 as coefficient.  After a
polynomial calculation is performed involving $\at$, the polynomial is
made square-free; that is, the exponent of each factor is reduced
to~1.  Given polynomial equation $\mathtt{P}(\at,x_1,\dots)$, its
square-free form is computed by (exactly) dividing $\mathtt{P}$ by the
GCD of $\mathtt{P}$ and its derivative with respect to $\at$:
\begin{equation}
{\mathtt{P}\left/\gcd\left(\mathtt{P},{\frac{\diff\mathtt{P}}{\diff\at}}\right)\right.}\label{normalize-at}
\end{equation}

Each of the remaining factors of the square-free polynomial equation
represents its degree of distinct solutions.  The square-free form
also removes factors not involving $\at$, as well as the ``content'',
the GCD of all the coefficients.  Expressions equivalent to 0 will
reduce to $\at=0$; therefore, this representation is a normal-form.

The existence of total orderings for variables and terms is well
established for Gr\"obner bases.  Using a total ordering, the
polynomial equation can be displayed in a deterministic way.  The
combination of a total ordering and this implicit normal form is thus
a canonical form for expressions.

The canonical form can be extended to equations not involving~$\at$.
In this case, the canonicalization procedure makes all factors
involving variables square-free.  The variable ordering governs which
high order term has a positive coefficient.

\section{Rational function differentiation}

Let
\begin{equation}
f(x)=\prod_{j\ne0} p_j(x)^j\label{f(x)}
\end{equation}
%%
be a rational function of $x$ where the primitive polynomials $p_j(x)$
are square-free and mutually relatively prime.

The derivative of $f(x)$ is
\begin{equation}
{\frac{\partial{f}}{ \partial{x}}}(x)=\sum_{j} j\,p_j(x)^{j-1}\,{p_j'}(x)\prod_{k\neq j}p_k(x)^k
  \label{f(x)'}
\end{equation}

\begin{lemma}\label{lemma:sum}
  Given square-free and relatively prime primitive polynomials $p_j(x)$,
  $\sum_{j}j\,{p_j'}(x)\prod_{k\neq j}p_k(x)$ has no factors in common
  with $p_j(x)$.
\end{lemma}

Assume that the sum has a common factor $p_h(x)$ such that $p_h(x)$
divides the sum:
\[
p_h(x)\left|\sum_{j}j\,{p_j'}(x)\prod_{k\neq j}p_k(x)\right.
\]

$p_h(x)$ divides all terms for $j\neq h$.  Because it divides the
whole sum, $p_h(x)$ must divide the remaining term $h\,p_h'(x)
\prod_{k\neq h} p_k(x)$.  From the given conditions, $p_h(x)$
does not divide $p_h'(x)$ because $p_h(x)$ is square-free; and $p_h(x)$
does not divide $p_k(x)$ for $k\neq h$ because they are relatively
prime.

\section{Rational function integration}

Separating square-and-higher factors from the sum in
equation~\eqref{f(x)'}:
\begin{equation}
{\frac{\partial{f}}{ \partial{x}}}(x) = \left[ \prod_j p_j(x)^{j-1} \right]
\left[\sum_j j\,{p_j'}(x)\prod_{k\neq j} p_k(x)\right]\label{f(x)'2}
\end{equation}

There are no common factors between the sum and product terms of
equation~\eqref{f(x)'2} because of the relatively prime condition of
equation~\eqref{f(x)} and because of Lemma~\ref{lemma:sum}.  Hence, this equation
cannot be reduced and is canonical.

Split equation \eqref{f(x)'2} into factors by the sign of the exponents,
giving:
\begin{equation}
{\frac{\partial{f}}{ \partial{x}}}(x)
 ={\frac{\prod_{j>2}p_j(x)^{j-1}}{\prod_{j<0} p_j(x)^{1-j}}}\,
 \overbrace{p_2(x)\,{\sum_{j}j\,{p_j'}(x)\prod_{k\neq j}p_k(x)}}^{\mathbf{L}}
\label{renum}
\end{equation}

The denominator is $\prod_{j\le0} p_j(x)^{1-j}$.  Its individual
$p_j(x)$ can be separated by square-free factorization.  The $p_j(x)$
for $j>2$ can also be separated by square-free factorization of the
numerator.  Neither $p_2(x)$ nor
${\sum_{j}j\,{p_j'}(x)\prod_{k\neq{j}}p_k(x)}$ have square factors; so
square-free factorization will not separate them.  Treating $p_2(x)$
as 1 lets its factor be absorbed into $p_1(x)$.  Note that $p_j(x)=1$
for factor exponents $j$ which don't occur in the factorization of
$\partial{f}/\partial{x}$.  All the $p_j(x)$ are now known except $p_1(x)$.  Once
$p_1(x)$ is known, $f(x)$ can be recovered by equation \eqref{f(x)}.
Let polynomial $\mathtt{L}$ be the result of dividing the numerator of
$\partial{f}/\partial{x}$ by $\prod_{j>2}p_j(x)^{j-1}$.
\begin{equation}
\overbrace{\sum_{j}j\,{p_j'}(x)\prod_{k\neq j} p_k(x)}^{\mathbf{L}}=
  \overbrace{\sum_{j\ne1}j\,{p_j'}(x)\prod_{1\neq k\neq j}p_k(x)}^{\mathbf{M}}\,p_1(x)
  +{p_1'}(x)\,\overbrace{\prod_{k\neq1}p_k(x)}^{\mathbf{N}}\label{p1(x)}
\end{equation}

Because they don't involve $p_1(x)$, polynomials $\mathtt{M}$ and
$\mathtt{N}$ in equation \eqref{p1(x)} can be computed from the
square-free factorizations of the numerator and denominator.  This
allows $p_1(x)$ to be constructed by a process resembling long
division.  The trick at each step is to construct a monomial $q(x)$
such that $\mathtt{M}\,q(x)+q'(x)\,\mathtt{N}$ cancels the highest term of
dividend $\mathtt{R}$ (which is initially $\mathtt{L}$).

Let $\deg(p)$ be the degree of $x$ in polynomial $p$.  Let
$\coeff(p,w)$ be the coefficient of the $x^w$ term of polynomial $p$
for non-negative integer $w$.

Note that $\deg(\mathtt{M})=\deg(\mathtt{N})-1$ because the
derivative of exactly one of the $p_j(x)$ occurs instead of $p_j(x)$
in each term of $\mathtt{M}$.  And
$\deg(q(x)\,\mathtt{M})=\deg(q'(x)\,\mathtt{N})$ because
$\deg(q'(x))=\deg(q(x))-1$.

The polynomial $p_1(x)$ can be constructed by the following procedure.
Let $\mathtt{A}$, $\mathtt{C}$, and $\mathtt{R}$ be rational expressions.
Only the numerators of $\mathtt{A}$ and $\mathtt{R}$ contain powers of
$x$.  Starting from polynomials $\mathtt{L}$, $\mathtt{M}$, and
$\mathtt{N}$:
\begin{verbatim}
A = 0
R = L
Nxd = deg(N)
while ( g = deg(num(R)) - Nxd + 1 ) >= 0:
    Rxd = deg(num(R))
    RxC = coeff(num(R),Rxd)
    C = RxC / ( coeff(M,Nxd-1) + g*coeff(N,Nxd) ) / denom(R)
    A = A + C * x^g
    R = R - C * ( M*x^g + N*diff(x^g,x) )
    if deg(num(R)) > Rxd:
        fail
    if 0 == R:
        return A
\end{verbatim}
At the end of this process, if $\mathtt{R}=0$, then $p_1(x)$ is the
numerator of $\mathtt{A}$; and the anti-derivative is
$f(x)=\prod_j{p_j(x)^j}/\mathrm{denom}\,(\mathtt{A})$.  Otherwise the
anti-derivative is not a rational function.

\medskip
Just as this algorithm works with $p_2(x)$ absorbed into $p_1(x)$, it
works with all of the $p_j(x)$ for $j>1$ absorbed into $p_1(x)$.  This
removes the need to factor the numerator and provides the opportunity
to enhance the algorithm to handle algebraic extensions.

\section{Algebraic extension}

Let variable $y$ represent one of the solutions of its defining
equation (reduction rule) represented by a polynomial $\mathtt{Y}=0$
which is irreducible over the integers.  For example $\mathtt{Y}$ would
be $y^3-x$ for a cube root of~$x$.

In order to normalize polynomials with regard to $\mathtt{Y}$, each
polynomial $\mathtt{P}$ containing $y$ is replaced by
$\mathtt{prem(P,Y)}$, the remainder of pseudo-division of $\mathtt{P}$ by
$\mathtt{Y}$, as described by Knuth~\cite{KnuthVol2}.

While that process normalizes polynomials, it doesn't normalize
ratios of polynomials, for instance:
\[
1/y^2=1/(\sqrt[3]{x})^2=\sqrt[3]{x}/x=y/x
\]

After the polynomials are normalized, if the denominator still
contains the extension $y$, it is possible to move $y$ to the
numerator by multiplying both numerator and denominator by the
$y$-conjugate of the denominator, then normalizing both numerator and
denominator by $\mathtt{Y}$.  The conjugate of a polynomial $\mathtt{P}$
with respect to $\mathtt{Y}$ can be computed by the following procedure
where $\deg(\mathtt{Q})$ is the degree of $y$ in polynomial~$\mathtt{Q}$ and
$\mathtt{pquo(Y,P)}$ and $\mathtt{prem(Y,P)}$ are the quotient and
remainder of pseudo-division of $\mathtt{Y}$ by~$\mathtt{P}$:
\begin{verbatim}
def conj(P):
    if deg(P) < deg(Y):
        Q = pquo(Y,P)
        R = prem(Y,P)
    else:
        Q = 1
        R = 0
    if deg(R) == 0:
        return Q
    else:
        return Q * conj(R)
\end{verbatim}

As discussed by Caviness and
Fateman~\cite{Caviness:1976:SRE:800205.806352}, multiple ring
extensions involving the same variable can be combined into a single
extension.  For the purposes of integration, combine the ring
extensions involving the variable of integration $x$ into a single
variable $y$ with its defining equation $\mathtt{Y}$.

\section{Rational function integration with algebraic extension}

With a single algebraic extension $y$ which is a function of
$x$, and the denominator free of $y$, and all the numerator factors in
$p_1(x,y)$, the previous development can be reformulated:
\begin{equation}
f(x,y)=\prod_{j\le1}p_j(x,y)^j\label{f(x,y)}
\end{equation}

The derivative of $f(x,y)$ with respect to $x$ is
\begin{equation}
{\frac{\partial{f}}{ \partial{x}}}(x,y)=\sum_{j\le1}j\,p_j(x,y)^{j-1}\,{p_j'}(x,y)\prod_{k\neq j}p_k(x,y)^k
  \label{f(x,y)'}
\end{equation}

Separating into numerator and denominator:
\begin{equation}
{\frac{\partial{f}}{ \partial{x}}}(x,y)
 ={\frac{\sum_{j}j\,{p_j'}(x,y)\prod_{k\neq j}p_k(x,y)}
        {\prod_{j\le0}p_j(x,y)^{1-j}}}
\label{renum2}
\end{equation}

This time, $\mathtt{L}$ is the whole numerator of
equation \eqref{renum2}.  Note that the denominator includes
$p_0(x,y)$; $p_0(x,y)$ does not contribute a term to $\mathtt{M}$
because its coefficient $j$ is 0.  Separating $p_1(x,y)$ from the
denominator factors:
\begin{equation}
\overbrace{\sum_{j}j\,{p_j'}(x,y)\prod_{k\neq j} p_k(x,y)}^{\mathbf{L}}=
  \overbrace{\sum_{j\le0}j\,{p_j'}(x,y)\prod_{k\neq j}p_k(x,y)}^{\mathbf{M}}\,p_1(x,y)
  +{p_1'}(x,y)\,\overbrace{\prod_{k\le0}p_k(x,y)}^{\mathbf{N}}\label{p1(x,y)}
\end{equation}

Because they don't involve $p_1(x,y)$, polynomials $\mathtt{M}$ and
$\mathtt{N}$ can be computed from the square-free factorization of the
denominator.  The trick at each step is to construct a polynomial
$t$ such that $\mathtt{M}\,t+t'\,\mathtt{N}$ cancels the
highest term of dividend $\mathtt{R}$ (initial $\mathtt{R}=\mathtt{L}$).

For polynomial $q$, $\deg(q'(x))=\deg(q(x))-1$ in the previous
section.  The derivatives of polynomials involving $x$ and its
algebraic extension $y$ are more complicated.  The derivative of
$y$ is found by differentiating the $y$ defining equation
$\mathtt{Y}=0$, then eliminating $y'$ from the chain rule for each term
of the polynomial it occurs in.  An example using the square of
$y=(x^4+a)^{1/3}$ demonstrates the reduction:
\[
{\frac{\diff{y^2}}{ \diff{x}}}=
  {\frac{8\,x^3}{3\,(x^4+a)^{1/3}}}=
  {\frac{8\,x^3\,(x^4+a)^{2/3}}{3\,(x^4+a)}}=
  {\frac{8\,x^3\,y^2}{3\,(x^4+a)}}
\]

The degree of $y$ does not change between $y^2$ and $\diff{y^2}/\diff{x}$.
However the degree of $x$ decreases; the degree of the denominator is
one more than the degree of the numerator of $\diff{y^2}/\diff{x}$.  This holds
for any algebraic extension defined by a primitive polynomial.

Let $\mathtt{A}$, $\mathtt{C}$, and $\mathtt{R}$ be rational expressions.
Let $\mathtt{Q}$ and $\mathtt{T}$ be polynomials of $x$ containing no
algebraic extensions of $x$.  Let $g=\deg_x\mathtt{R}-\deg_x\mathtt{N}+1$.
The addition of 1 is to compensate for the reduction in the degree of
$x$ in $p_1'(x,y)$.

When there is no algebraic extension, $t=x^g$.  If there is an
algebraic extension $y$, let $q$ be the denominator of normalized
${\diff{y}/\diff{x}}$, $i$ be the integer quotient ${g/\deg_x{q}}$, and set
$g$ to the remainder of ${g/\deg_x{q}}$.  Then:
\[
t=q^i\,x^g\,y^h
\]

%% Because $y$ is algebraic function of $x$, $\deg(\diff{y}/\diff{x},x)=\deg(y,x)-1$.

The polynomial $p_1(x,y)$ can be constructed by the following
procedure given polynomials $\mathtt{L}$, $\mathtt{M}$, and $\mathtt{N}$:
\begin{verbatim}
A = 0
R = L
Q = denom( normalize( diff(y,x) ) )
Nyd = deg(N,y)
NyC = coeff(N,y,Nyd)
Nxd = deg(NyC,x)
while 0 < 1:
    Ryd = deg(num(R),y)
    RyC = coeff(num(R),y,Ryd)
    Rxd = deg(RyC,x)
    h = Ryd - Nyd
    g = ( Rxd - Nxd + 1 )
    if 0 = deg(Q,x):
        T = x^g
    else:
        i = quotient(g,deg(Q,x))
        g = remainder(g,deg(Q,x))
        T = Q^i * x^g * y^h
    B = normalize( N*diff(T,x) + M*T )
    C = coeff(RyC,x,Rxd) * denom(B) / denom(R) /
        coeff(coeff(num(B),y,Ryd),x,Rxd)
    A = A + C * T
    R = R - C * B
    if 0 = R:
        return A
    if deg(num(R),y) > Ryd:
        fail
    if deg(num(R),y) == Ryd and
       deg(coeff(num(R),y,deg(num(R),y)),x) >= Rxd:
        fail
\end{verbatim}
The looping continues only as long as the degree of $\mathtt{R}$
decreases.  If this process succeeds, then the numerator of $\mathtt{A}$
is $p_1(x,y)$; and the anti-derivative is
$f(x,y)=\mathtt{A}\,\prod_{j<1}p_j(x,y)^j$.

\section{Transcendental extension}

Some transcendental extensions involving a variable can be
described by separable first-order differential equations:
\begin{equation}
{{\log(x)}'}={\frac{1}{ x}}\,{x'}\label{diff-log}
\end{equation}

Where $a$ is independent of $x$, applying the chain-rule:
\begin{equation}
{{\log(x^a)}'}={\frac{a\,x^{a-1}}{ x^a}}\,{x'}={\frac{a}{ x}}\,{x'}\label{diff-log-a}
\end{equation}

Eliminating $x'/x$ from equations \eqref{diff-log} and \eqref{diff-log-a}
results in:
\[
{{\log(x^a)}'}=a\,{{\log(x)}'}
\]

Taking the anti-derivative of both sides:
\begin{equation}
\log(x^a)=a\,\log(x)+C\label{log(x^a)}
\end{equation}

Thus, $\log(x)$ is sufficient to be the canonical extension for both
$\log(x)$ and $\log(x^a)$.  Logarithm is multi-valued in the complex
plane.  Like radical functions, $\log(x)$ assumes one branch.  The
canonical form described here produces equations which are true for
any branch of $\log$ assuming all references to $\log(x)$ are the same
branch.  Note that the branch $\log(x)$ is not linked to the
branch of $\log(y)$.
%% Note the difference from algebraic extensions; the branches for
%% $\sqrt{x}$ and $\sqrt{y}$ aren't linked, while the branches for
%% $\log{x}$ and $\log{y}$ are.

To describe the inverse function of $\log(x)$, substitute $y$ for
$\log(x)$ and $\exp(y)$ for $x$:
\begin{equation}
{\frac{{y}'}{ {\exp(y)}'}}={\frac{1}{\exp(y)}}\qquad{\frac{{\exp(y)}'}{ {y}'}}=\exp(y)\label{diff-exp}
\end{equation}

Part of normalization for an expression (or equation) reduces the
expression modulo the defining rules of the extensions appearing in
that expression.  For a transcendental extension, this is its
defining differential equation.

These reductions serve to normalize expressions (involving variables)
with unnested extensions.  While (irreducible) algebraic expressions
involving unnested transcendental extensions are canonical, when a
transcendental function is composed with its inverse, an algebraic
expression (without transcendental extensions) can result.
%% The rest of this section addresses those nested transcendental
%% expressions which reduce to algebraic forms.

\section{Reduction of nested transcendental extensions}

For $\log(\exp(y))$, the reduction results from integrating the
combination of defining differential equations~\eqref{diff-log}
and~\eqref{diff-exp} with $x=\exp(y)$:
%% $${\log(x)'\over x'}={1\over x}\qquad {\exp(y)'\over y'}={\exp(y)}\label{exp-log-defs}$$
\[
{\log(\exp(y))'}={\frac{\exp(y)'}{\exp(y)}}=y'
 \qquad \log(\exp(y))=y+C
\]

For $\exp(\log(x))$, the defining differential equations~\eqref{diff-log}
and~\eqref{diff-exp} combined with $y=\log(x)$ is separated, integrated
($\ln[=\log]$ is the result of the integration), then exponentiated:
\[
{\frac{\exp(\log(x))'}{\exp(\log(x))}}={\frac{x'}{ x}}
 \qquad{\ln(\exp(\log(x)))}=\ln(x)
 \qquad \exp(\log(x))=x
\]

The defining equations for arc-tangent and tangent are:
\begin{equation}
\arctan(x)'={\frac{x'}{ x^2+1}}\qquad
 {\frac{\tan(\theta)'}{\theta'}}=\tan(\theta)^2+1\label{tan-arctan-defs}
\end{equation}
\[
x=\tan(\theta)\qquad
 \arctan(\tan(\theta))'={\frac{\tan(\theta)'}{\tan(\theta)^2+1}}=\theta'
 \qquad \arctan(\tan(\theta))=\theta+C
\]
\[
\theta=\arctan(x)\qquad
 {\frac{\tan(\arctan(x))'}{\arctan(x)'}}=\tan(\arctan(x))^2+1
 \qquad{\frac{\tan(\arctan(x))'}{\tan(\arctan(x))^2+1}}={\frac{x'}{ x^2+1}}
\]

For both $\tan(\arctan(x))$ and $\exp(\log(x))$, composition through
the defining equations results in a form with polynomial function
$\Phi$:
\begin{equation}
{\frac{y'}{\Phi(y)}}={\frac{x'}{\Phi(x)}}\label{Phi-form}
\end{equation}

Equation \eqref{Phi-form} is a stronger constraint than $y'=x'$ in
that it implies $y=x$ without a constant of integration.
Equation \eqref{Phi-form} is not the only form reducing to an
algebraic relationship between $y$ and $x$.
Equation \eqref{Phi-b-form} where $b$ is a nonzero integer may also
reduce to an algebraic relation:
\begin{equation}
{\frac{y'}{\Phi(y)}}={\frac{b\,x'}{\Phi(x)}}\label{Phi-b-form}
\end{equation}

Separating $y(x)$ into a ratio of relatively prime polynomials
$y(x)=f(x)/g(x)$:
\begin{equation}
{\frac{(f/g)'}{\Phi(f/g)}}=
   {\frac{f'\,g-f\,g'}{ g^2\,\Phi(f/g)}}={\frac{b\,x'}{\Phi(x)}}\qquad
   ({f'\,g-f\,g'})\,{\Phi(x)}=b\,g^2\,\Phi(f/g)\,x'\label{y=f/g}
\end{equation}

$f$, $f'$, $g$, $g'$, $x$, $x'$, and $\Phi(x)$ are all polynomials.
Thus, $({f'\,g-f\,g'})\,{\Phi(x)}$ is a polynomial.  If
equation \eqref{y=f/g} is to hold, then $g^2\,\Phi(f/g)\,x'$ must also
be a polynomial, which will only be the case when the denominator of
$\Phi(f/g)$ has degree 1 or 2 times the degree of $g$.  When
$g^2\,\Phi(f/g)\,x'$ does not equal $({f'\,g-f\,g'})\,{\Phi(x)}$, the
scaled composition cannot be reduced to an algebraic relation.

For both $\Phi(y)=y$ and $\Phi(y)=y^2\pm1$, it is the case that
$g^2\,\Phi(f_b/g_b)={\Phi(x)}^b$, and $\Phi$ is related to
$y_b=f_b/g_b$ with $b\ge2$:
\[
\left({f_b'\,g_b-f_b\,g_b'}\right)\,{\Phi(x)}=b\,{\Phi(x)}^b\,x'\quad\to\quad
  {f_b'\,g_b-f_b\,g_b'}=b\,{\Phi(x)}^{b-1}\,x'
\]

In the case of $\Phi(y)=y$,
\begin{equation}
y_1=x\qquad y_{b}=y_{b-1}\,x\qquad {\frac{y'}{ y}}={\frac{b\,x'}{ x}}\to y=x^b\label{log-y_b}
\end{equation}
%% $${a\,y'\over y^2}={b\,x'\over x^2}\to y^a=x^b\qquad$$
%% $${a\,(1+x^2)\,y'}={b\,(1+y^2)\,x'}$$

In the case of $\Phi(y)=1+y^2$, using the sum-of-angles formula (a la
Chebyshev):
\begin{equation}
\tan(b\,\theta)={\frac{\tan((b-1)\,\theta)+\tan(\theta)}{1-\tan((b-1)\,\theta)\,\tan(\theta)}}\label{Chebyshev}
\end{equation}
%% ={\sum_{k\ \rm{odd}}(-1)^{(k-1)/2}\,{b\choose k}\,\tan^k\theta\over\sum_{k\ \rm{even}}(-1)^{k/2}\,{b\choose k}\,\tan^k\theta}

Let $\theta=\arctan x$, $y_{1}=x$, $y_b=f_b/g_b$:
\begin{equation}
y_{b}={\frac{y_{b-1}+x}{1-y_{b-1}\,x}}=
  {\frac{f_{b-1}+g_{b-1}\,x}{ g_{b-1}-f_{b-1}\,x}}\qquad
  y_{-b}={\frac{y_{-b+1}-x}{1+y_{-b+1}\,x}}=
  {\frac{g_{-b+1}\,x-f_{-b+1}}{ g_{-b+1}+f_{-b+1}\,x}}\label{arctan-y_b}
\end{equation}
%% ={(1-x^2)\,y_{b-2}+2\,x\over(1-x^2)-y_{b-2}\,2\,x}
%% ={(1-x^2)\,f_{b-2}+g_{b-2}\,2\,x\over(1-x^2)\,g_{b-2}-f_{b-2}\,2\,x}
%% ={\sum_{k\ \rm{odd}}(-1)^{(k-1)/2}\,{b\choose k}\,x^k\over\sum_{k\ \rm{even}}(-1)^{k/2}\,{b\choose k}\,x^k}

Because $f_{1}=x$ and $g_{1}=1$, the difference of the degrees of
$f_{b}$ and $g_{b}$ alternate between $1$ and $-1$ with the parity of
$b$.

Extensions for irreducible $\Phi$ polynomials will always reduce
to identity when directly composed with their inverse function; those
which can reduce scaled composition to more complicated algebraic
expressions are limited to degrees of 1 or 2 by
equation \eqref{y=f/g}.  Any $y_{b}$ recurrence must be symmetrical in
$x$ and $y_{b-1}$.  There are few symmetrical candidates for $y_{b}$
which might have transcendental-to-algebraic reductions.  They were
checked by formula \eqref{y=f/g} for $y_2(x)=f_2/g_2$.

For hyperbolic tangent $y_{b}=\tanh(b\,\atanh x)$, $\Phi(y)=1-y^2$, $y_{1}=x$, $y=f/g$ and:
\begin{equation}
y_{b}={\frac{y_{b-1}+x}{1+y_{b-1}\,x}}=
  {\frac{f_{b-1}+g_{b-1}\,x}{ g_{b-1}+f_{b-1}\,x}}\qquad
  y_{-b}={\frac{y_{-b+1}-x}{1-y_{-b+1}\,x}}=
  {\frac{g_{-b+1}\,x-f_{-b+1}}{ g_{-b+1}-f_{-b+1}\,x}}\label{arctanh-y_b}
\end{equation}
%%   {(1+x^2)\,y_{b-2}+2\,x\over(1+x^2)+y_{b-2}\,2\,x}=
%%   {(1+x^2)\,f_{b-2}+g_{b-2}\,2\,x\over(1+x^2)\,g_{b-2}+f_{b-2}\,2\,x}
%% $$\tanh(n\,\theta)={\tanh((n-1)\,\theta)+\tanh(\theta)\over1+\tanh((n-1)\,\theta)\,\tanh(\theta)}$$

But $\Phi(x)=1-x^2$ is reducible, having factors $1-x$ and $1+x$; so
it is not associated with a lone canonical extension.  Instead
the integral uses extensions $\log(x+1)$ and $\log(x-1)$:
\[
\int{\frac{\diff{x}}{1-x^2}}={\frac{\log(x+1)-\log(x-1)}{2}}
\]

Equation \eqref{Phi-form} is symmetrical.  Both sides can be scaled by
nonzero integers:
\begin{equation}
{\frac{a\,y'}{\Phi(y)}}={\frac{b\,x'}{\Phi(x)}}\label{Phi-a-b-form}
\end{equation}

A composition of defining differential equations yielding a
form \eqref{Phi-a-b-form} results in a relation $y_{b}=x_{a}$.  For
$\Phi(t)=t$, $y=(\sqrt[a]{x})^b$.  For other $\Phi$, the result is
a polynomial relation $y_{b}=x_{a}$.

Because ${x'/\Phi(x)}+{x'/\Phi(x)}=2\,{x'/\Phi(x)}$, using
equations~\eqref{log-y_b} and~\eqref{arctan-y_b} without recurrence
directs how to compose with a sum of inverse transcendental functions.
\begin{equation}
\begin{aligned}
  \exp(\log x_1+\log x_2)&\to y=x_1\,x_2\\
  \tan(\arctan x_1+\arctan x_2)&\to y=(x_1+x_2)/(1-x_1\,x_2)
\end{aligned}
\label{composition-of-sum}
\end{equation}

A related problem is to normalize $\log(x\,z)\to\log x+\log z$.
Taking the total differential of $\log(x\,z)$ yields:
  \[
\log(x\,z)'={\frac{x'\,z+x\,z'}{ x\,z}}={\frac{x'}{ x}}+{\frac{z'}{ z}}
\]
which is separable and integrable to $\log x+\log z$.  This procedure
works for $\arctan()$ as well:
  \[
\int\arctan\left({\frac{x+z}{1-x\,z}}\right)'=\int{\frac{x'}{1+x^2}}+\int{\frac{z'}{1+z^2}}
  =\arctan x+\arctan z
\]

These example nested transcendental functions simplify to rational
functions (polynomial ratios):
\[
\begin{aligned}
\exp(3\,\log x)&\to{\frac{y'}{ y}}={\frac{3\,x'}{ x}}\to y_3=x^3\\
\tan(4\,\arctan x)&\to
  {\frac{y'}{1+y^2}}={\frac{4\,x'}{1+x^2}}\to
  y_4={\frac{{4\,x-4\,x^{3}}}{{1-6\,x^{2}+x^{4}}}}\\
\tan(7\,\arctan x)&\to
  {\frac{y'}{1+y^2}}={\frac{7\,x'}{1+x^2}}\to
  y_7={\frac{{-7\,x+35\,x^{3}-21\,x^{5}+x^{7}}}{{-1+21\,x^{2}-35\,x^{4}+7\,x^{6}}}}
\end{aligned}
\]

\newpage

\section{Nested transcendental extensions}

%% The upper incomplete gamma function $\Gamma(s,x)$ is notated here as
%% $\Gamma_{\!s}(x)$.  With non-integer argument $s$:
%% $${\partial{\Gamma_{\!s}(x)}\over\partial{x}}={-x^{s-1}\over\exp{x}}$$

%% These equations must hold with change-of-variable.  Substituting
%% $x^2+a$ for $x$:
%% $$\int{2\,x\,\diff{x}\over\left(x^2+a\right)^2+1}=\arctan\left(x^2+a\right)$$

%% \proclaim Lemma 2. {
%% Given differentiable functions $f(x)$ and $g(y)$, where $y=f(x)$ and
%% $${\partial g/\partial y\over g(y)}=
%%   {\partial f/\partial x\over f(x)}$$}\par

%% $$\tan(2\,\theta)={2\,\tan(\theta)\over1-\tan^2(\theta)}\qquad
%%   \tan{\theta\over2}={-1\pm\sqrt{1+\tan\theta}\over\tan\theta}$$

%% Let $g(f)$ and $f(x)$ be transcendental functions.  Consider the
%% derivative of their composition:
%% $$g(f(x))'=g'(f(x))\,f'(x)\,x'\qquad{g(f)'\over g'(f)}={f'(x)\,x'}\label{composition-derivative}$$
%% %% $$g(a\,f(x))'=g'(a\,f(x))\,a\,f'(x)\,x'\qquad{g(a\,f(x))'\over g'(a\,f(x))}={a\,f'(x)\,x'}$$

%% Thus, their derivative always forms a separable differential equation.

\medskip

%% Inverse function theorem:
%% $$x=f(y(x))\qquad 1={f(y(x))'\over{x'}}=f'(y)\,y'(x)\qquad {f'(y)}={1\over y'(x)}\qquad f(y)=\int{1\over y'(x)}\,x'$$
%% $$(x\,\exp{x})'=[1+x]\,\exp{x}\,x'$$

The Lambert $\W$ function is an inverse of $y\,\exp y$.  Its
derivative is used as its defining differential equation:
\begin{equation}
{\frac{\partial\W(z)}{\partial{z}}}={\frac{\W(z)}{ z\;[1+\W(z)]}}\label{W}
\end{equation}
\[
{\frac{w'}{{z'}}}={\frac{ w}{ z\;[1+ w]}}\qquad
  \left[1+{\frac{1}{{w}}}\right]\;{w'}={\frac{{z'}}{{z}}}\qquad
  w+{\ln{w}}={\ln{z}}\qquad z=w\,\exp{w}
\]
%% $$\qquad z=w\,\exp{w}\qquad{{z'}}=\left[1+{w}\right]\;\exp{w}\;w'$$

The derivative of the composition of $\W$ with its inverse,
$x\,\exp{x}$, reduces to a separable differential equation:
\begin{equation}
{\frac{\W(x\,\exp{x})+1}{\W(x\,\exp{x})}}\;{\W(x\,\exp{x})}'={\frac{x+1}{ x}}\,x'\label{W-1}
\end{equation}

Equation \eqref{W-1} is of the form:
\begin{equation}
{\frac{\W'}{\Phi(\W)}}={\frac{x'}{\Phi(x)}}\qquad\Phi(y)={\frac{y}{{y+1}}}\label{sep-W}
\end{equation}

Thus, $\W(x\,\exp{x})=x$.  Because $\Phi(y)$ is not a polynomial,
reductions from forms like equation~\eqref{Phi-a-b-form} are not
possible for~$\W$.
Differentiating ${\W(x)\,\exp(\W(x))}$, then reducing by
equations~\eqref{diff-exp} and~\eqref{W}:
\begin{equation}
\begin{aligned}
  {(\W(y)\;\exp(\W(y)))'}&={\exp(\W(y))\;\W(y)'}+{\exp(\W(y))'\;\W(y)}\\
  &={\frac{\exp(\W(y))\;[1+\W(y)]\;\W(y)\;y'}{ y\;[1+\W(y)]}}
\end{aligned}
\label{sep-W1}
\end{equation}
%% $$\eqalign{{(\W(y)\;\exp(\W(y)))'}&={\exp(\W(y))\;\W(y)'}+{\exp(\W(y))'\;\W(y)}\cr
%%   &={\exp(\W(y))\;\W(y)'}+{\exp(\W(y))\,\W(y)\,\W(y)'}\cr
%%   &=\exp(\W(y))\,\left[1+{\W(y)}\right]\,\W(y)'\cr
%%   &=\exp(\W(y))\,\left[1+{\W(y)}\right]\,{\W(y)\over y\,[1+\W(y)]}y'\cr
%%   &={\exp(\W(y))\;[1+\W(y)]\;\W(y)\;y'\over y\;[1+\W(y)]}\cr}$$
%% &={\exp(\W(y))\;[1+\W(y)]\;\W(y)'}\cr
\begin{equation}
{\frac{(\W(y)\;\exp(\W(y)))'}{\W(y)\;\exp(\W(y))}}={\frac{y'}{ y}}\label{sep-W2}
\end{equation}

Thus, ${\W(x)\,\exp(\W(x))=x}$.
%% There is a non-differential recurrence for $\W$, but it is doubly
%% recursive and won't reduce to purely algebraic expressions:
%% $$\W(x)+\W(y)=\W\left(x\,y\left[{1\over\W(x)}+{1\over\W(y)}\right]\right)$$

\section{Transcendental constants}

This system is canonical for algebraic extensions involving variables.
It cannot be canonical for algebraic constants because the many roots
of unity make alternate factorizations possible, and the system
requires that polynomials be a unique factorization domain.  Algebraic
constants can be represented; but not all possible reductions will
happen as part of the normalization processes.

This is also the reason that trigonometric functions were not encoded
with $\log()$ and $\exp()$ using imaginary transcendental constants
($\pi\,\mathrm{i}$).

Transcendental constants can be represented by an uninstantiated
differential equation and a rational argument value:
\[
\exp1=\mathrm{e}\qquad\arctan1=\pi/4
\]

If a transcendental function is composed with the transcendental
constant, then their differential equations are combined.  If the
resulting equation is separable, then the recurrence is used to
construct the algebraic result as before.  $\exp0=1$ and $\arctan0=0$
are deducible from the recurrence without explicit training.

\section{Rational function integration with transcendental extension}

For an integer power $w$ of an algebraic extension $y(x)$:
\[
\deg_y{\frac{\partial{y^w}}{ \partial{x}}}=w\qquad\deg_x{\frac{\partial{y^w}}{ \partial{x}}}=-1
\]

Transcendental extensions behave differently:
\medskip

\noindent
\begin{tabular}{@{}p{\thirdwidth}p{\thirdwidth}p{\thirdwidth}@{}}
  \qquad$y=\exp(x^w)$    & \qquad$\deg_y{\frac{\partial{y}}{ \partial{x}}}=1$
                         & \qquad$\deg_x{\frac{\partial{y}}{ \partial{x}}}=w-1$\\
  \qquad$y=\log^w(x)$    & \qquad$\deg_y{\frac{\partial{y}}{ \partial{x}}}=w-1$
                         & \qquad$\deg_x{\frac{\partial{y}}{ \partial{x}}}=-1$\\
\end{tabular}

\refstepcounter{table}\label{transcendental extensions}%
{\centering\textbf{Table \thetable\quad transcendental extensions}\par}

%% There are two ways a variable representing a transcendental
%% extension can arise in integration, from a non-constant polynomial
%% in $p_0$, or from appearing in any of the $p_j$ polynomials.

\bibliographystyle{unsrt}
\bibliography{ratint}

\end{document}

% LocalWords:  tex Caviness coeff atanh erf fre Ritt GCD obner Nxd Nyd
% LocalWords:  prem pquo factorizations renum num denom Fateman deg
% LocalWords:  Rxd RxC NyC Ryd RyC unnested exponentiated arctan defs
% LocalWords:  Chebyshev arctanh sep uninstantiated ratint Jaffer
% LocalWords:  monomial endverbatim canonicalization
